\documentclass[11pt,aps,preprint,nofootinbib,superscriptaddress,tightenlines]{revtex4-2}
\usepackage[T1]{fontenc}
\usepackage{amsmath,amssymb,mathtools}
\let\lambdabar\relax   % revtex defines it; newpxmath does too
\usepackage{newpxtext,newpxmath}
\usepackage{bm}
\usepackage[final]{microtype}
\usepackage{xcolor}
\usepackage{booktabs}
\usepackage{dsfont}
\newcommand{\one}{\mathds{1}}
\definecolor{darkblue}{RGB}{20,60,140}
\usepackage[colorlinks=true,linkcolor=darkblue,citecolor=darkblue,urlcolor=darkblue]{hyperref}

\def\st{\begin{equation}}
\def\stp{\end{equation}}
\def\Eq#1{eq.~(\ref{#1})}
\def\Sect#1{Sect.~\ref{#1}}
\def\Tab#1{Table~\ref{#1}}
\let\oldparagraph\paragraph
\renewcommand{\paragraph}{\par\bigskip\oldparagraph}
\newcommand{\ask}{\textbf{Ask}}

\begin{document}
\sloppy

\title{Focusing the code: plan for the pruning and the restructure}
\author{D.~Teaney}
\affiliation{Department of Physics and Astronomy, Stony Brook University,
Stony Brook, New York 11794, USA}
\date{\today}

\begin{abstract}
The code has accumulated options: every change of default left the old path reachable through a setting.  This plan reduces it to its
mission, deterministic runs in two dimensions and stochastic runs in two dimensions, with one choice for each part of the method: the
number volume with the energy Hamiltonian, the generalized leapfrog, the inverse-moment viscous step, and the Liouville measure.  It then
makes the inputs strict, rewrites the comments, and restructures the main driver.  Almost every stage must leave the numbers bit-identical; the
few that change them on purpose are separate stages with their own checks.  The plan is written for the session at NERSC.  The decisions
behind it were taken in a review with the owner on 2026-11-27; the discussion is in \nolinkurl{docs/code_review_notes.md}.
\end{abstract}

\maketitle
\tableofcontents

%=============================================================================
\section{What stays}
\label{sec:stays}

After the work the code does three kinds of run, all with periodic boxes in two dimensions:
\begin{itemize}
\item deterministic runs (the ideal step and the deterministic viscous step), with the physical dimension $d=2$ by default or $d=3$ (a fluid
  independent of $z$, simulated in the plane);
\item stochastic runs, $d=2$;
\item hybrid Monte Carlo, for the start of the stochastic runs and for equilibrium measurements.
\end{itemize}
Each part of the method has one form:

\begin{center}
\small
\begin{tabular}{p{0.16\textwidth}p{0.36\textwidth}p{0.38\textwidth}}
\toprule
part & stays & goes \\
\midrule
Hamiltonian, volume & energy Hamiltonian, $V_a=1/n_a$ (the n-scheme) & smoothed Hamiltonian, entropy volume \\
ideal step & generalized leapfrog & drift--kick--drift leapfrog, the grid step \\
viscous step & \texttt{:moment} & \texttt{:particle}, \texttt{:common}, \texttt{:bond}, \texttt{:two\_level} \\
measure & Liouville & Einstein \\
representation & the particle variables only & the change of representation, the gradient back map \\
stochastic step & single vote, $\lambda=1$; Hastings on or off & votes \texttt{:bonds}, \texttt{:even}; $\lambda=\infty$; \texttt{residual}, \texttt{shift} \\
dimension & simulation in 2D; $d=2$ or 3 & the 1D path \\
start & \texttt{cold}, \texttt{hot}, \texttt{hmc} & \texttt{thermal}, \texttt{SR\_BATH=totals}, Poisson positions \\
renormalization & tree level, split, S-scheme & (unchanged; split may stay the default) \\
\bottomrule
\end{tabular}
\end{center}

Kept although little used: the kernels Wendland C4 and C6 (undeveloped; to be brought alive after this work, \Sect{sec:after}), the jitter of
the lattice (with its own random numbers, Q10), the frozen lattice of the stochastic step, the flag \texttt{SPH\_HASTINGS} (and with it the
coefficients stored in \texttt{Star}), the generic equation of state behind \texttt{AbstractEOS}, \texttt{energy\_pressure}, the test switches
\texttt{SR\_NOVISC}, \texttt{SR\_NOHAM}, \texttt{SR\_DTFAC}, \texttt{SR\_RECORD\_E}, \texttt{SR\_CS2}, \texttt{SR\_ZETAS}, and
\nolinkurl{examples/legacy/}.  The \texttt{Particle} struct is not changed in this plan.

%=============================================================================
\section{For the NERSC session}
\label{sec:nersc}

The rules of \nolinkurl{CLAUDE.md} apply.  In addition:
\begin{enumerate}
\item \emph{Read first.}  \nolinkurl{docs/code_review_notes.md} (the review, with file and line references), \nolinkurl{docs/codeplan.tex} (the
  test suite and the run record), and \nolinkurl{nersc/LOG.md}.
\item \emph{Branch.}  Work on a new branch \texttt{focus} from \texttt{main}, in a new worktree.  One commit per stage, pushed.  \texttt{main} is
  merged only on the owner's word.
\item \emph{The gate.}  Every stage that is not marked otherwise passes
  \st
  \texttt{test/run.sh all -r pre-prune}
  \stp
  with the regression bit-identical.  A difference in the last bits is not accepted silently: find whether the arithmetic changed or only the
  order of operations (for instance the compiler's fusion of multiply--add, which depends on the surrounding code; several functions are
  \texttt{@noinline} for this reason), and report it.
\item \emph{Removed settings stop.}  Every setting removed in a stage goes into \texttt{RETIRED\_SETTINGS} (\nolinkurl{src/settings.jl}) in the
  same stage, with a message that names the tag \texttt{pre-prune} for reruns of old campaign lines.  Removed values of a setting that stays (for
  instance \texttt{SPH\_RENORM=1}) stop in the same way.
\item \emph{Scripts and tests.}  In each stage, a script in \nolinkurl{examples/} or a test that only exercises a removed path is deleted; one that
  also does something that stays is changed to the path that stays.  Campaign files are not edited: they are the record of what was run.
\item \emph{Running campaigns.}  The stages change the checkpoint format (Q14 at the latest).  Before the branch is merged, check that no
  campaign is still running or may need to resume from a checkpoint of the old code.
\item \emph{The S-scheme work} is paused until after this plan (owner, 2026-11-27): the kernels C4, C6 and the non-conformal run of
  the S-scheme come after the restructure (\Sect{sec:after}).  \texttt{focus} runs without merges from \texttt{main}.
\item \emph{Budget.}  Tests and smoke runs in \texttt{-q debug} or an interactive node.  The only campaign of this plan is the check of Q9;
  \ask{} with its estimate.
\item \emph{Check-ins.}  Stop at the points marked \ask{} and whenever a gate fails.  Record each stage in \Sect{sec:results} (what was removed,
  the commit, the gate's result).
\end{enumerate}

\paragraph{Decisions of the owner for this session (2026-11-27).}  They replace the \ask{} points they name.
\begin{enumerate}
\item \emph{Budget:} up to 150 node-hours for this plan without asking; estimates and job IDs in \nolinkurl{nersc/LOG.md}.
\item \emph{Q11:} as written in Q11 below.
\item \emph{Q9:} no stop for the cost; the HMC comparison campaign runs within the budget.  Stop only if a check fails (\texttt{hot} or
  \texttt{cold} not bit-identical to \texttt{pre-prune}, or the HMC statistics disagree).
\item \emph{Q0:} if the only regression differences are the known ones of the regression's comment (the record of $H$ since
  \texttt{3fa7a3e}, the nominal-bath start since \texttt{2bc03dd}), record them in \Sect{sec:results} and continue to Q1.  Stop if anything
  else differs.
\item \emph{S-scheme work in parallel:} none; \texttt{focus} has the files to itself.
\item \emph{Still asked:} the new record layout in Q14 before committing it, the final merge of \texttt{focus} into \texttt{main} (Q15), and
  whenever a gate fails.
\end{enumerate}

%=============================================================================
\section{The stages}
\label{sec:stages}

The order removes the leaves first: the options of the stochastic step, then the paths that only the removed options use, then the larger
alternatives.  Stages Q1--Q8 are bit-identical against \texttt{pre-prune}; Q9 and Q10 change numbers on purpose; Q11 changes behavior only where
a step fails; Q12--Q14 are bit-identical against the end of Q11.

\subsection*{Q0. The baseline}
\begin{itemize}
\item Tag the head of \texttt{main} as \texttt{pre-prune} and push the tag.
\item \texttt{test/run.sh all -r main} at the tag, on an interactive node.  Record in \Sect{sec:results} which checks pass and which known
  differences remain (the regression's comment lists them).
\item The regression of \nolinkurl{test/run.sh} loops over \texttt{SPH\_HAMILTONIAN} (energy, smoothed) and runs the equilibrium driver with
  \texttt{SPH\_RENORM=0 SR\_BATH=totals}.  Add the defaults (split, the nominal bath, the HMC start) next to the tree level, so that the gate
  covers what production runs use.  The \texttt{smoothed} entries are removed in Q7 and the \texttt{totals} ones in Q9.
\end{itemize}
\ask{} the owner with the list of known differences before Q1.

\subsection*{Q1. The options of the stochastic step}
Remove, in \nolinkurl{src/stochastic.jl}:
\begin{itemize}
\item $\lambda$ and its path: \texttt{resid\_deterministic!}, \texttt{moment\_fit}, \texttt{GUARD\_SUBDIV}, the $\lambda$ argument of the
  proposals and votes, \texttt{SR\_LAMBDA} in the drivers.  The residual move always has $\lambda=1$.
\item The votes \texttt{:bonds} and \texttt{:even}: the shares $\sigma_{ab}$ and the even split in \texttt{ask\_star!},
  \texttt{GUARD\_EVENSPLIT}, the \texttt{vote} argument and field.  The single vote stays.  Note the commit \texttt{pre-prune} in
  \nolinkurl{docs/stochastic.tex} where the bond vote is described: the owner may come back to it.
\item The fields \texttt{residual} and \texttt{shift} of \texttt{StochOptions} (always true).
\end{itemize}
Stays: \texttt{SPH\_HASTINGS} and the path without Hastings (single vote), the frozen lattice.  Users to change:
\nolinkurl{examples/stochastic_moving.jl}, \nolinkurl{test/teststochastic_guard.jl}.

\subsection*{Q2. The measure}
Liouville only.  Remove the field \texttt{measure} of \texttt{SimParams}, \texttt{default\_measure}, \texttt{SPH\_MEASURE}, the argument
\texttt{liouville} of \texttt{vote\_entropy}, the votes and the sweeps, the Einstein branches (\texttt{vote\_entropy} returning $S_a$ alone, the
$\Sigma\ln T_a$ term of the acceptance of \nolinkurl{src/hmc.jl}, the end-of-step branch of \texttt{stochastic\_step!}), and the checks ``needs
the Liouville measure'' in \nolinkurl{src/renorm_split.jl} and \nolinkurl{src/sscheme.jl}.

\emph{Care.}  \texttt{HMCTarget} builds its Gaussian proposal with \texttt{particle\_moments(..., liouville = false)}, the Einstein weight.
Changing it changes the proposals and hence the numbers of every HMC start.  In this stage keep that weight for the proposal (an internal
function, named for what it is), so the stage stays bit-identical; the proposal moves to the Liouville weight in Q9.  The keyword defaults
\texttt{liouville = false} of \nolinkurl{src/measure.jl} go.

\subsection*{Q3. The change of representation}
Remove the gradient back map (\texttt{gradient\_freeze}, \texttt{gradient\_backmap!}, the field \texttt{gradient} of \texttt{StochOptions},
\texttt{GUARD\_BACKMAP}), the gradient-corrected conversion map of \nolinkurl{src/pointwise.jl} (\texttt{kernel\_star!}, \texttt{kernel\_star},
\texttt{star\_fit}, \texttt{inverse\_susceptibility}, \texttt{gradient\_energy}) and \texttt{pointwise\_state}, \texttt{pointwise\_beta!} if
nothing else uses them, and the deterministic viscous steps that belong to the change of representation in \nolinkurl{src/viscous.jl}
(\texttt{viscous\_step\_first!}, \texttt{viscous\_step\_exact!}, \texttt{viscous\_step\_grad!}: check against \texttt{viscous\_step\_energy!},
which stays).  Delete \nolinkurl{test/testchangerep.jl}.  The two particle maps of \nolinkurl{src/pointwise.jl} stay.

\subsection*{Q4. The one-dimensional path}
Remove \texttt{WendlandC2\_1D}, the branch \texttt{sim\_dim == 1} of \texttt{SimParams} and its keyword, the function \texttt{sim\_dim} (or make it
the constant 2), the 1D cases of the \texttt{default\_*} functions and of \texttt{check\_stochastic\_dims}, \texttt{make\_1d\_particles!} and its
branch in \texttt{setup\_ic}, \texttt{bond\_weights\_1d} and the \texttt{oneD} branch of the bond dissipation, the 1D constants of
\texttt{viscous\_rate\_bound}, the one-row case of \texttt{neighbor\_range}.  Delete \nolinkurl{examples/testsound_viscous_1d.jl},
\nolinkurl{examples/testgaussian_1d.jl}, the 1D branch of \nolinkurl{examples/stochastic_thermalize.jl}, the 1D part of
\nolinkurl{test/testdims.jl}.  The field \texttt{phys\_dim} stays (default 2; $d=3$ for deterministic runs).

\subsection*{Q5. The viscous step}
\texttt{:moment} only.  Remove \texttt{:particle}, \texttt{:common}, \texttt{:bond} (and the name \texttt{:changerep}), the field
\texttt{dissipation} and \texttt{default\_dissipation}, \texttt{SPH\_DISSIPATION}, \texttt{bond\_weights\_comoving} and the $\Lambda_a$
correction if \texttt{:moment} does not use them, \texttt{is\_common}, \texttt{lab\_weights}, the residual factor's \texttt{:common} branch, and
\texttt{bond\_h} (field, keyword, \texttt{bond\_weight}); \texttt{viscous\_rate\_bound} then uses $h$.  \texttt{:moment} currently requires the
energy Hamiltonian: after Q7 that check goes too.  Delete \nolinkurl{test/testcommonkernel.jl} or move what it checks of \texttt{:moment} into
\nolinkurl{test/testmoment.jl}.  \nolinkurl{test/testmoment.jl} compares \texttt{:moment} with \texttt{:common} (its item 4, reported only): drop
that item.

\subsection*{Q6. The ideal step}
The generalized leapfrog only.  Remove the drift--kick--drift body of \texttt{leapfrog\_step!} (which then calls \texttt{generalized\_step!}), the
grid step (\nolinkurl{src/grid.jl}, \texttt{record\_grid!}, the \texttt{GRID} branches and series of \nolinkurl{examples/equilibrium.jl}, the
checks ``not with the grid step''), the field \texttt{ideal}, \texttt{default\_ideal}, \texttt{SPH\_IDEAL}.  Delete \nolinkurl{examples/grid_*.jl}
and \nolinkurl{test/testgridfft.jl}.

\subsection*{Q7. The Hamiltonian and the volume}
The n-scheme only.  Remove the smoothed Hamiltonian (the bodies of \texttt{smooth\_to\_particles!}, \texttt{update\_velocities!},
\texttt{hamiltonian\_force!} in \nolinkurl{src/ideal.jl}; each then is its \texttt{*\_energy!} version), the entropy volume (the
\texttt{number = false} branches), the fields \texttt{hamiltonian} and \texttt{volume}, \texttt{default\_hamiltonian}, \texttt{SPH\_HAMILTONIAN},
\texttt{SPH\_VOLUME}, and the \texttt{smoothed} loop of the regression.  Delete \nolinkurl{src/legacy/} (the two-level form and the old sweep;
the tag keeps them) and point to the tag in \nolinkurl{docs/codeplan.tex}.  Keep the round trip $S_{\rm smooth}=s_a n_a$, $V_a=s_a/S_{\rm
smooth}$: removing it changes the last bits (\Sect{sec:after}).  The \texttt{Particle} struct is not changed.

\subsection*{Q8. The rest}
Remove \texttt{make\_uniform\_particles} (and the square lattice if nothing else uses it), \texttt{smooth\_to\_grid}, and the Poisson positions
of the driver (\texttt{SR\_POSITIONS=poisson}; \nolinkurl{test/testmoment.jl} builds its own Poisson stars).  Check \nolinkurl{test/run.sh load}:
every script that remains still loads.

\subsection*{Q9. The start of a run (numbers change)}
\label{sec:q9}

\paragraph{The target state.}  One function,
\begin{multline}
\texttt{target\_state}(s_{\rm rel,0}, v, \text{eos}) \to \big(T_{\rm rel,0},\, v,\, \gamma,\, e_0,\, w_0,\, S=\gamma s_{\rm rel,0},\\
M_0=w_0\gamma^2 v,\, \beta^0=\gamma/T_{\rm rel,0},\, \bm\beta=\gamma\bm v/T_{\rm rel,0}\big),
\end{multline}
from the equation of state alone.  The set-up, the draw, hybrid Monte Carlo, the normalization of KE and the record all read it.
\texttt{SR\_BATH} goes, with \texttt{beta\_from\_averages} and the set-up-particle values of $w$, $T$ in the driver; \texttt{moments\_table}
stays.

\paragraph{One flag for the start.}  \texttt{SR\_START} $\in\{$\texttt{cold}, \texttt{hot}, \texttt{hmc}$\}$, default \texttt{hmc}, replaces
\texttt{SR\_START}, \texttt{SR\_POSITIONS} and \texttt{SR\_BATH}:
\begin{itemize}
\item \texttt{cold}: the lattice as set up, no fluctuations; the run thermalizes.  For demonstrations.
\item \texttt{hot}: the lattice, $(e_a,\bm m_a)$ drawn at the target $\beta^\mu$, then shifted so that the total energy and momentum equal those of
  the cold start.  Cold and hot are then on the same shell and must reach the same equilibrium.  For demonstrations.
\item \texttt{hmc}: the draw, then hybrid Monte Carlo at the target $\beta^\mu$ (positions, and $(e_a,\bm m_a)$ through the heat bath).  No
  redraw after it, no shift.  For production.  A debugging flag \texttt{SR\_HMC\_SHIFT\_TOTALS} (Bool, default false) shifts the final sample
  to the cold totals.  The HMC settings (\texttt{SR\_HMC\_*}) are read only here.
\end{itemize}
The cold totals $(E_{\rm cold},\bm P_{\rm cold})=(\sum_a e_a,\sum_a \bm m_a)$ of the lattice are computed once, right after the set-up and before
any draw, and kept in the run state and the record.  \texttt{SR\_START=thermal} (\texttt{thermal\_init!}) and \texttt{SR\_START=eq} stop with a
message naming the new values.  The extra redraw after HMC (today \texttt{initial\_state!} with \texttt{:eq}) goes.  The proposal of the HMC heat
bath moves to the Liouville weight (Q2).

\paragraph{One loop of hybrid Monte Carlo.}  The loop of HMC updates with random trajectory lengths is written in five drivers
(\nolinkurl{examples/equilibrium.jl}, \nolinkurl{eos_measure.jl}, \nolinkurl{hmc_positions.jl}, \nolinkurl{moment_fluctuations.jl},
\nolinkurl{volume_fluctuations.jl}).  Move it into \nolinkurl{src/hmc.jl} as \texttt{hmc\_equilibrate!} and use it everywhere.

\paragraph{Checks.}  (i) \texttt{hot} is bit-identical to \texttt{pre-prune} with \texttt{SR\_POSITIONS=lattice SR\_START=eq}.  (ii)
\texttt{cold} is bit-identical to \texttt{pre-prune} with \texttt{SR\_START=cold SR\_POSITIONS=lattice}.  (iii) \texttt{hmc} is a different,
statistically equivalent start: compare with \texttt{pre-prune} at $s_a=54$ and 16, at rest and $v_0=0.8$, the equal-time quantities (KE/eq,
the static spectra of $M$ and $S$, $\langle H\rangle$) and the transverse and longitudinal correlators.  \ask{} with the estimate before
submitting.  (iv) \texttt{SR\_HMC\_SHIFT\_TOTALS=true} reproduces the cold totals to rounding.

\subsection*{Q10. The random numbers of the jitter (numbers change)}
The jitter of \texttt{make\_hex\_particles!} uses the global \texttt{RNG} of \nolinkurl{src/sph.jl}.  Give it its own generator,
\texttt{Xoshiro(params.seed)}, and remove the global \texttt{RNG} (\nolinkurl{test/testneighbor.jl} uses it once).  Jittered positions change; runs
without jitter are bit-identical.  The tests that use jitter (\nolinkurl{test/testmoment.jl}, \nolinkurl{test/testhastings.jl},
\nolinkurl{test/testdissipative.jl}) must pass with the new positions.

\subsection*{Q11. A failed step of the leapfrog}
\texttt{generalized\_leapfrog!} returns $-1$ when it reaches a state without a Hamiltonian; hybrid Monte Carlo rejects the trajectory, but
\texttt{generalized\_step!} ignores the return, so a run continues from that state.

\emph{Decision (owner, 2026-11-27).}  On a failed step of the generalized leapfrog (a $-1$ return, or the implicit iterations not
converging), restore $x$ and $m$ and redo the step as two half steps, recursively, down to $\Delta t/16$.  Count each retry in a new
guard counter, reported by the monitor and in the record's \texttt{[outcome]}.  If the step still fails at $\Delta t/16$, save the state at
the start of the step and the failing particles for diagnosis, mark the record failed with the reason, and stop with a nonzero exit.  The
run must never continue from a state without a Hamiltonian.  Bit-identical when no step fails.

\subsection*{Q12. Strict inputs}
\begin{itemize}
\item \emph{A table of the settings} in \nolinkurl{src/settings.jl}: every valid name, its default and a one-line meaning.  Each
  \texttt{setting(...)} call takes its default from the table.  Defaults that depend on other choices are written as such in the table.
\item \emph{Unknown names stop.}  At load, any \texttt{SR\_*} or \texttt{SPH\_*} in the environment that is not in the table stops the run, with
  the nearest valid name in the message.  Retired names keep their own message.
\item \emph{Bad values stop.}  Bool: \texttt{1}/\texttt{true}/\texttt{yes} or \texttt{0}/\texttt{false}/\texttt{no} (any case), anything else
  stops.  Symbol: checked against its allowed values.  Numbers: as now.
\item \emph{Unused names are recorded.}  Just before the main loop, the names set in the environment but never read go into an
  \texttt{[unused]} table of the run record, with their values.  Nothing is printed and the run does not stop.
\item \emph{A listing.}  \texttt{julia src/settings.jl --list} prints the table.
\item Optional: \nolinkurl{nersc/check_campaign.jl}, which checks the names and values of every line of a campaign file before \texttt{sbatch}.
\end{itemize}
Numbers bit-identical; the record gains \texttt{[unused]}.

\subsection*{Q13. The comments}
Rewrite the comments of \nolinkurl{src/} and of the drivers that stay.  Short lines; one point per paragraph; a blank comment line (\texttt{\#}
alone) between logically distinct points.  The dated history of options that no longer exist goes; where a choice was made for a reason that
still matters (for instance \texttt{@noinline} for bit-identity), the reason stays.  Gate: the regression is bit-identical, and the diff of the
stage with comments stripped is empty.

\subsection*{Q14. The main driver}
Restructure \nolinkurl{examples/equilibrium.jl}:
\begin{itemize}
\item \emph{\texttt{RunConfig}}: every input of a run, read once through the table of Q12, grouped (the physics, the numerics, the start, the
  output, the test switches), printed at the start and written to the record.  The command-line arguments go into it too.
\item \emph{\texttt{RunState}}: a declared struct in place of the named tuple \texttt{st} (the system, the random numbers, the time, the
  references, the recorders and series).  This is what a checkpoint holds.
\item \emph{The parts}: \texttt{one\_realization} split into the set-up, the start (Q9), the evolution, the recording, and the checkpoint and
  output.  What other drivers need (the set-up of the fluid at the target state, the start, the Strang loop with a recording callback, the
  checkpoints) moves into \nolinkurl{src/}, so that \nolinkurl{examples/stability.jl}, \nolinkurl{eos_measure.jl} and the others use the same
  pieces.
\item \emph{Compatibility}: checkpoints of the old code do not resume (say so in \nolinkurl{nersc/README.md}).  The keys of the record that
  \nolinkurl{nersc/collect_runs.jl} and \nolinkurl{examples/boost_analysis.jl} read are kept or those scripts are changed in the same stage.
  \ask{} the owner with the new layout of the record before committing it.
\end{itemize}
Gate: bit-identical data against the end of Q13 for the same settings.

\subsection*{Q15. The documents}
\nolinkurl{CLAUDE.md}, \nolinkurl{README.md}, \nolinkurl{GETTING_STARTED.md}, the settings table of \nolinkurl{nersc/README.md}, and
\nolinkurl{docs/codeplan.tex}: describe what stays, and point to \texttt{pre-prune} for what went.  \ask{} the owner to read them, then to merge
\texttt{focus} into \texttt{main}.

%=============================================================================
\section{After this plan}
\label{sec:after}

Not part of the plan, recorded so that they are not lost:
\begin{itemize}
\item The kernels C4 and C6: tabulate the positional counterterm for them, so that they run with \texttt{SPH\_POSITIONAL=1}.  Then the comparison of the kernels (C2 against C4).
\item The non-conformal run of the S-scheme: split with $\omega$, the measure form at $b_\rho=b$ and at $b_\rho=D/2+1$ (about 5--7 node-hours;
  \nolinkurl{docs/sscheme_plan.tex}, sec:beta), after the choice of $b_\rho$ and of the reference state off the conformal point, and a pass of
  the tests of the S-scheme and of $\omega$ at $c_s^2=1/3$.
\item Split stays (owner, 2026-11-27): it may remain the default.
\item The round trip $S_{\rm smooth}=s_a n_a$ and the fields of \texttt{Particle} that the removed paths used (\texttt{M\_smooth} as a smoothed
  field, \texttt{Teff}, \texttt{F\_over\_S}, \texttt{dotx}, $\Delta s$, $\Delta m$): a later change of the struct, not bit-identical.
\item The tuning of hybrid Monte Carlo (\texttt{nmh}, \texttt{widen}, the tolerances) into the record.
\item The seeds: a list of the blocks of \texttt{SR\_SEEDOFFSET} in use, in \nolinkurl{nersc/LOG.md}.
\end{itemize}

%=============================================================================
\section{Results}
\label{sec:results}

One entry per stage: what was removed or changed, the commit, the result of the gate, and any decision taken.

The gate is \texttt{test/run.sh all -r pre-prune} (Q0--Q11) or \texttt{-r focus-q11} (Q12--Q14), on an interactive or a regular node.
Two changes of the harness apply to every stage: \nolinkurl{test/regress/step.jl} is the present version on both sides (the owner agreed,
2026-11-27), and \nolinkurl{test/regress/h5cmp.jl} does not compare the \texttt{hast\_lq\_*} datasets, the statistics of the Hastings
log ratios summed with atomic adds by the threads: their last bits differ between two runs of identical code (the Q0 baseline showed it).
They are reported.

\paragraph{Q0} (tag \texttt{pre-prune} = \texttt{f5237d9}; \texttt{1d69ec2}, \texttt{32da36c}).  The regression runs the drivers at the
defaults too.  Baseline (interactive 62102226, the code of the tag against the tag): unit, physics, load pass; \texttt{step.jl} and the
stability driver bit-identical; the equilibrium driver bit-identical but for \texttt{hast\_lq\_*}; the unit logs differ only in a wall
time of hybrid Monte Carlo (now masked).  The known differences of the regression's comment (the record of $H$, the nominal bath)
belong to references before the merge and do not appear against the tag.

\paragraph{Q1} (\texttt{60a127e}, \texttt{a61e506}).  Gate 62103015: the drivers bit-identical (but for \texttt{hast\_lq\_*}); one unit
failure, \nolinkurl{test/testmoment.jl} calling \texttt{propose\_resid!} with the removed $\lambda$ (fixed in Q5).

\paragraph{Q2} (\texttt{1911977}, \texttt{81b6837}).  Gate (on 62103015): drivers bit-identical; \texttt{testmoment} as in Q1 and
\texttt{testeqstart}, whose quadrature integrated the Einstein target (fixed in \texttt{81b6837}).  \emph{Deviation from the plan}: with
the Einstein measure gone, the smoothed Hamiltonian cannot run the stochastic step or the equilibrium start (it had that measure), so
the regression's smoothed entries went in Q2, not Q7.  The proposal of the HMC heat bath kept the flat weight (\texttt{flat\_weight})
until Q9, as planned.

\paragraph{Q3--Q8} (\texttt{15aeb7d}, \texttt{2843434}, \texttt{10a4084}, \texttt{e6d2074}, \texttt{68d29c3}, \texttt{4d6e300}).  Gated
together at the head of Q8 (62103676): unit, physics, load pass; the regression bit-identical to \texttt{pre-prune} (\texttt{step.jl};
the equilibrium driver at the tree level and at the defaults, split with the HMC start, at rest and $v_0=0.6$; the stability driver);
the unit logs differ only where the tests changed with the removed options (and the count of settings, a wall time), and the deleted
tests are reported as only in \texttt{pre-prune}.  Decisions: the one-row case of \texttt{neighbor\_range} stays (a 2D box less than $4h$ high
needs it); \nolinkurl{test/testdissipative.jl} moved to \texttt{:moment} (the continuum $K$ of linear fields, the exact fit, the finite
differences: all hold to roundoff), \nolinkurl{test/testcommonkernel.jl}, \nolinkurl{test/testparticle_sweep.jl},
\nolinkurl{test/testkappa.jl}, \nolinkurl{test/testgridfft.jl}, \nolinkurl{test/testchangerep.jl} deleted with their paths;
\nolinkurl{examples/moment_det.jl} keeps \texttt{:moment} alone, \nolinkurl{examples/moment_gaussian.jl} and
\nolinkurl{examples/commonkernel_waves.jl} (comparisons with \texttt{:common}) deleted.

\paragraph{Q9} (\texttt{cfd6988}).  Checks (i), (ii): \texttt{hot} and \texttt{cold} bit-identical to \texttt{pre-prune} with
\texttt{SR\_POSITIONS=lattice SR\_START=eq} and \texttt{=cold}, at the tree level and with split, at rest and $v_0=0.8$.  Check (iv):
\texttt{SR\_HMC\_SHIFT\_TOTALS} reproduces the cold totals to $10^{-15}$ (\nolinkurl{test/teststart.jl}, at $s_a\approx540$; at $s_a\approx16$
in motion the uniform shift can push a particle past the light cone, and \texttt{shift\_to\_totals!} stops).  Check (iii): campaign
\texttt{q9\_start} (62104058, 62104059; \nolinkurl{nersc/campaign_q9_start.txt}; 16 seeds per state, $t=100$, split with $\omega$),
compared by \nolinkurl{examples/q9_start_compare.jl}: KE/eq, $\langle H\rangle$, the static spectra of $M_L$, $M_T$, $S$ at the five
smallest $k$, and $C_{T,L}(\tau)/C_{T,L}(0)$ at $\tau=0.5,1,2$, at $s_a=16$, 54, $v_0=0$, 0.8: 92 comparisons, the rms of
(new$-$old)$/\sigma$ is 1.00, 26 beyond $1\sigma$, 7 beyond $2\sigma$, none beyond $3\sigma$ (largest 2.56).  KE/eq: $1.040(5)\to1.039(5)$,
$1.008(3)\to1.008(3)$ at rest.  The new start is statistically equivalent: check (iii) passes.  The stability driver now starts \texttt{hot} at the
target (before: the totals solve), so its regression against \texttt{pre-prune} differs from Q9 on.

\paragraph{Q10} (\texttt{5aa40c9}).  \texttt{testmoment}, \texttt{testhastings}, \texttt{testdissipative} pass with the new jittered
positions; \texttt{testneighbor} keeps its numbers with its own generator.

\paragraph{Q11} (\texttt{504ea20}; tag \texttt{focus-q11}).  The owner's decision as written.  \nolinkurl{test/testleapfrog_fail.jl}: a
step 1000 times the acoustic one fails at every depth: four retries, a \texttt{LeapfrogFailure} with the state at the start of the
step, the driver exits 1 with the record marked failed and the state saved.  Gate 62104423 (the Q11 head against \texttt{pre-prune}): unit,
physics, load pass; \texttt{step.jl} and the equilibrium driver (tree and default, hot and cold, rest and boost) bit-identical; the
stability driver differs (its start, Q9); the unit-test output differs only in the tests changed or removed (as at Q8).

\paragraph{Q12} (\texttt{9576776}).  The table \texttt{SETTING\_SPECS}; the batch and test variables are in it as \emph{infra} (not read
through \texttt{setting}, not listed as unused).  Defaults that depend on the driver (\texttt{SR\_KHMAX}, \texttt{SR\_START},
\texttt{SR\_SEED0}, \ldots) are marked so and given by the driver; every other call takes the table's.  The optional
\nolinkurl{nersc/check_campaign.jl} is not written.  Gate 62104585 against \texttt{focus-q11}: unit, physics, load pass; \texttt{step.jl} and the stability
driver bit-identical, the unit-test output identical.  Its driver comparisons failed in the harness (the detection of a ref before Q9
matched a comment of \nolinkurl{src/settings.jl}); fixed in \nolinkurl{test/run.sh}, and covered by the Q15 gate.

\paragraph{Q13} (\texttt{3c28f88}).  The comments of the 18 files of \nolinkurl{src/} rewritten; \nolinkurl{test/regress/astcmp.jl}
(the parsed code, line numbers stripped) identical to the Q12 head.

\paragraph{Q14} (\texttt{6425dc8}).  \nolinkurl{src/run.jl} (\texttt{setup\_fluid}, \texttt{start\_run!}, \texttt{RunState},
\texttt{resume\_run}, \texttt{strang\_loop!}); \nolinkurl{examples/equilibrium.jl} on \texttt{RunConfig} with a \texttt{[config]}
section in the record (the owner waived the question of its layout).  A checkpoint is a \texttt{RunState}; one of the code before Q14
stops with a message.  The driver bit-identical to \texttt{focus-q11} in 12 configurations; \texttt{testcheckpoint},
\texttt{testleapfrog\_fail}, \texttt{testrunrecord}, \texttt{teststart} pass.

\paragraph{Q15} (\texttt{40fad27}).  \nolinkurl{CLAUDE.md}, \nolinkurl{README.md}, \nolinkurl{GETTING_STARTED.md},
\nolinkurl{nersc/README.md}, \nolinkurl{docs/codeplan.tex} rewritten for the one scheme.  Gate 62105280 against \texttt{focus-q11}:
unit, physics, load pass; \texttt{step.jl}, the equilibrium driver (tree and default; hot, cold and hmc; rest and boost) and the
stability driver bit-identical; the unit-test output differs only by the printed configuration (\texttt{testcheckpoint}) and the
new file \nolinkurl{src/run.jl} in the load list (\texttt{testrunrecord}).  The plan is done; the merge into \texttt{main} waits for
the owner.

\end{document}
